\documentclass[t, aspectratio=169]{beamer}
\usepackage{fontspec}
\IfFontExistsTF{Times New Roman}{\setmainfont{Times New Roman}}{\setmainfont{Liberation Serif}}
\usepackage{amsmath}
\usepackage{amssymb}
\usepackage{booktabs}
\usepackage{graphicx}
\usepackage{xcolor}
\usepackage{array}
\usepackage{tikz}
\usetikzlibrary{shapes.callouts, positioning, arrows.meta, fit, backgrounds}

\definecolor{hustblue}{RGB}{0,60,113}
\definecolor{goodgreen}{RGB}{30,120,60}
\definecolor{warnorange}{RGB}{200,90,20}
\definecolor{badred}{RGB}{170,30,30}

\def\slideContentLeftMargin{0.5cm}
\def\slideContentRightMargin{0.5cm}
\def\hustHeaderHeight{3.0cm}
\def\slideContentBotMargin{2.8cm}

\usepackage[blue, 169]{beamerthemeHUST}
\usefonttheme{serif}
\renewcommand{\familydefault}{\rmdefault}

\setbeamercolor{frametitle}{fg=hustblue}
\setbeamerfont{frametitle}{size=\fontsize{15}{17}\selectfont, series=\bfseries}
\setbeamertemplate{frametitle}{
  \nointerlineskip
  \begin{tikzpicture}[remember picture, overlay]
    \node[anchor=north west, xshift=0.35cm, yshift=-0.95cm, text width=12.3cm, align=left]
    at (current page.north west)
    {\usebeamerfont{frametitle}\usebeamercolor[fg]{frametitle}\insertframetitle};
  \end{tikzpicture}
}
\setbeamertemplate{footline}{\vbox to 1.8cm{}}

\renewcommand{\titlePageTitleX}{-0.5cm}
\renewcommand{\titlePageTitleY}{-0.5cm}
\renewcommand{\hustTitleAlign}{\alignLeft}
\renewcommand{\hustSubtitleAlign}{\alignLeft}
\renewcommand{\hustAuthorAlign}{\alignLeft}
\renewcommand{\hustInstAlign}{\alignLeft}

\renewcommand{\hustsectionpage}[1]{%
  {%
    \setbeamertemplate{background}{\includegraphics[width=\paperwidth,height=\paperheight]{\hustbg{section}}}%
    \begin{frame}[plain,c]
      \centering \Huge \textcolor{white}{\textbf{#1}}
    \end{frame}
  }%
}

\graphicspath{{images/}}

\newcommand{\bignum}[4]{%
  \begin{minipage}[t]{#4}
    \centering
    {\fontsize{32}{36}\selectfont\bfseries\color{#3} #1}\\[2pt]
    {\small #2}
  \end{minipage}%
}
\newcommand{\bignumS}[4]{%
  \begin{minipage}[t]{#4}
    \centering
    {\fontsize{22}{26}\selectfont\bfseries\color{#3} #1}\\[1pt]
    {\scriptsize #2}
  \end{minipage}%
}

\title{\large Báo cáo cuối kỳ}
\subtitle{\normalsize Tích hợp dữ liệu khí tượng đa nguồn Đông Nam Á}
\author{Ngô Văn Quang \quad $\cdot$ \quad Nguyễn Cao Lương \quad $\cdot$ \quad Trương Tuấn Nghĩa}
\institute{SoICT, Đại học Bách khoa Hà Nội \\
IT5427 -- Tích hợp và xử lý dữ liệu lớn -- Nhóm 16\\[3pt]
{\small Giảng viên: TS.\ Vũ Tuyết Trinh}}

\begin{document}

\husttitlepage

\begin{frame}[shrink=5]{Nội dung}
\begin{columns}[T]
\begin{column}{0.5\textwidth}
\begin{enumerate}
\item \hyperlink{sec1}{Tổng kết tiến độ 160 tác vụ}
\item \hyperlink{sec2}{Ingestion -- TV1}
\item \hyperlink{sec3}{Schema \& Warehousing -- TV2}
\end{enumerate}
\end{column}
\begin{column}{0.5\textwidth}
\begin{enumerate}
\setcounter{enumi}{3}
\item \hyperlink{sec4}{Spark, Streaming \& PINN -- TV3}
\item \hyperlink{sec4b}{Phát hiện dị thường khí hậu}
\item \hyperlink{sec5}{Dashboard -- TV4}
\item \hyperlink{sec6}{Rủi ro, giới hạn \& kết luận}
\end{enumerate}
\end{column}
\end{columns}
\end{frame}

\resetPageCount
\showPageNumber
\useStyleHustFull
\useThinBar

%=============================================================
\section{Tổng kết tiến độ}
%=============================================================
\hypertarget{sec1}{}
\hustsectionpage{1. Tổng kết tiến độ\\160 tác vụ}

\begin{frame}{Pipeline tích hợp 5 nguồn khí tượng đa dạng thành 1 kho dữ liệu thống nhất, có thể truy vấn được}
\begin{center}
\bignum{9.348.470}{bản ghi thật\\5 nguồn}{goodgreen}{4.6cm}
\bignum{119}{vị trí\\(113 trạm + 6 điểm)}{goodgreen}{3.4cm}
\bignum{9}{nước Đông Nam Á}{goodgreen}{3.4cm}
\end{center}
\vspace{0.3cm}
\centering\small
GHCN-Daily, NOAA ISD, ERA5 (ARCO), Open-Meteo, NASA POWER -- hội tụ qua Mediated Schema
\texttt{WeatherObs}, giải quyết cả 2 loại dị biệt cốt lõi: cấu trúc/ngữ nghĩa (PIVOT/PARSE) và
nhận dạng/toàn vẹn (Survivorship + Entity Resolution). Dữ liệu đủ phạm vi đề cương: ERA5 5 năm
(2020--2024, 9 biến, theo giờ), ISD 10 năm, GHCN toàn lịch sử 1943--2026 -- \textasciitilde 23\,GB dữ liệu thô.
\end{frame}

\begin{frame}{Mỗi trạng thái tác vụ được chấm theo đúng tiêu chí Kiểm thử đã đặt ra từ đề cương, không tự nhận suông}
\centering\small
\begin{tabular}{@{}lccc@{}}
\toprule
Tuần & \textcolor{goodgreen}{\textbf{Xong}} & \textcolor{warnorange}{\textbf{Một phần}} & \textcolor{badred}{\textbf{Chưa làm}} \\
\midrule
Tuần 1 (Đề cương) & 30/40 & 4/40 & 6/40 \\
Tuần 2 (Tiến độ) & 30/40 & 5/40 & 5/40 \\
Tuần 3 (Hội tụ) & 19/40 & 12/40 & 9/40 \\
\midrule
\textbf{Tổng 120 tác vụ đã rà soát} & \textbf{79} & \textbf{21} & \textbf{20} \\
\bottomrule
\end{tabular}
\vspace{0.3cm}

\small Tuần 4 (Hoàn thiện \& Bảo vệ) chủ yếu là hành động con người (diễn tập, trình bày, nộp
bài) -- phần kỹ thuật tương ứng (freeze, Q\&A, tài liệu cuối) đã làm xong song song.
\end{frame}

%=============================================================
\section{Ingestion}
%=============================================================
\hypertarget{sec2}{}
\hustsectionpage{2. Ingestion\\(TV1)}

\begin{frame}{IQR phát hiện đủ 5/5 outlier cực đoan trong khi Z-score chỉ bắt được 1/5 -- một phát hiện thống kê thật}
\begin{center}
\resizebox{!}{3.4cm}{%
\begin{tikzpicture}[remember picture,
  bar/.style={fill=badred!60}, barok/.style={fill=goodgreen!70}]
\draw[-{Latex[length=2.5mm]}, thick] (0,0) -- (0,3.2) node[above] {\scriptsize tỷ lệ phát hiện};
\draw[-{Latex[length=2.5mm]}, thick] (0,0) -- (7,0);
\draw[bar] (1,0) rectangle (2.5,0.6) node[midway, above, font=\small] {};
\node[font=\small] at (1.75,0.85) {\textbf{20\%}};
\node[font=\scriptsize] at (1.75,-0.35) {Z-score};
\draw[barok] (4,0) rectangle (5.5,3.0);
\node[font=\small] at (4.75,3.25) {\textbf{100\%}};
\node[font=\scriptsize] at (4.75,-0.35) {IQR};
\node[rectangle callout, callout absolute pointer={(1.75,0.8)}, fill=badred!10, text=badred,
  rounded corners, inner sep=5pt, align=center, font=\scriptsize]
  at (2.8,2.3) {"Masking effect":\\outlier tự làm\\std phình to};
\end{tikzpicture}%
}
\end{center}
\vspace{0.1cm}
\centering\small
Tiêm 5 outlier cực đoan vào dữ liệu N(25,2): std bị kéo từ 2,0 lên 72,2, khiến ngưỡng 3$\sigma$ quá
rộng để bắt chính các outlier đó -- IQR (dựa trên tứ phân vị) không bị ảnh hưởng. Đã chọn IQR làm
chỉ số chính cho Data Quality Report.
\end{frame}

\begin{frame}{Đã tải đủ dữ liệu theo phạm vi đề cương: kho dữ liệu tăng 8,7 lần lên 9,3 triệu bản ghi}
\centering\footnotesize
\begin{tabular}{@{}lrrl@{}}
\toprule
Nguồn & Trước & Hiện tại & Phạm vi hiện tại \\
\midrule
ERA5 (nội suy vào 113 trạm) & 8.136 & \textbf{4.954.824} & 2020--2024, 9 biến, theo giờ \\
NOAA ISD & 358.273 & \textbf{3.464.430} & 26 trạm, 2015--2024 \\
GHCN-Daily & 655.164 & 655.166 & 41 trạm, 1943--2026 \\
Open-Meteo + NASA POWER & 54.900 & \textbf{274.050} & 6 điểm, 2020--2024 \\
\midrule
\textbf{Tổng} & 1.076.473 & \textbf{9.348.470} & 119 vị trí, 109 partition \\
\bottomrule
\end{tabular}
\vspace{0.15cm}

\footnotesize Kiểm thử ERA5 ``không thiếu tháng nào'': 540/540 file, đủ 43.848/43.848 giờ mỗi biến, 0\% NaN.
Nút thắt là \textbf{băng thông}, không phải số luồng: mỗi chunk ARCO-ERA5 là 1 giờ \emph{toàn cầu}, nên
lấy vùng ĐNA vẫn phải tải \textasciitilde 1\,TB (\textasciitilde 24 giờ) để giữ lại 22\,GB. Script tải theo tháng,
chạy lại được sau khi bị ngắt.
\end{frame}

\begin{frame}{Kiểm tra lại trên dữ liệu đầy đủ phát hiện và sửa 4 lỗi số liệu thật trong pipeline}
\centering\scriptsize
\begin{tabular}{@{}p{3.4cm}p{4.9cm}p{4.9cm}@{}}
\toprule
Lỗi & Nguyên nhân & Hệ quả / sửa \\
\midrule
Gió Open-Meteo \textbf{phóng đại 3,6 lần} & API trả km/h (\texttt{hourly\_units}) nhưng ghi thẳng vào cột \texttt{wind\_speed\_ms} & Sai mọi so sánh gió đa nguồn -- parser đọc đơn vị và quy đổi, có test \\ \addlinespace[2pt]
\textbf{Tương quan giả} (gió--độ ẩm +0,57) & Thay NaN bằng 0 trước khi tính; mỗi nguồn thiếu biến khác nhau & Tính từng cặp: lộ ra nhiệt độ--độ ẩm \textbf{$-$0,60}, nhiệt độ--áp suất \textbf{$-$0,61} (đúng vật lý) \\ \addlinespace[2pt]
Benchmark partition so với \textbf{bản sao cũ} & Bản không partition chỉ ghi 1 lần; luôn in ``nhanh hơn'' & Đo lại công bằng: partition chậm hơn 1,5--2,4 lần ở 9,3 triệu dòng \\ \addlinespace[2pt]
Bức xạ ERA5 \textbf{sai đại lượng} & Mapping dùng bức xạ đỉnh khí quyển (TOA) & Đổi sang bức xạ bề mặt, cùng đại lượng với NASA POWER / Open-Meteo \\ \addlinespace[2pt]
\bottomrule
\end{tabular}
\end{frame}

%=============================================================
\section{Schema \& Warehousing}
%=============================================================
\hypertarget{sec3}{}
\hustsectionpage{3. Schema \&\\Warehousing (TV2)}

\begin{frame}{Entity Resolution theo khoảng cách toạ độ khớp được 22/26 trạm ISD với GHCN, giải quyết gap đã biết từ Tuần 2}
\begin{center}
\bignumS{22/26}{cặp GHCN$\leftrightarrow$ISD\\khớp được (84,6\%)}{goodgreen}{4.5cm}
\bignumS{0,53 km}{khoảng cách\\trung bình}{goodgreen}{4.5cm}
\end{center}
\vspace{0.2cm}
\centering\small
\texttt{warehouse/entity\_resolution.py} -- Haversine nearest-neighbor, ngưỡng 5km. Nhiều cặp khớp
tại \textbf{0,000km} (cùng tọa độ WMO gốc) -- xác nhận đây đúng là 1 trạm vật lý được 2 cơ quan gán
2 mã độc lập, không phải trùng hợp. 4 trạm ISD còn lại không có GHCN tương ứng trong ngưỡng.
\end{frame}

\begin{frame}{Đối chiếu 108 nghìn cặp trạm--ERA5 xác nhận ERA5 nội suy tin cậy cho nhiệt độ, gió nhưng KHÔNG cho mưa}
\centering\small
\begin{tabular}{@{}lrrl@{}}
\toprule
Biến & RMSE & MAE & Đánh giá \\
\midrule
\texttt{temp\_c} & 1,14°C & 0,85°C & \textcolor{goodgreen}{Chấp nhận được} \\
\texttt{wind\_speed\_ms} & 1,02 m/s & 0,75 m/s & \textcolor{goodgreen}{Chấp nhận được} \\
\texttt{precip\_mm} & 19,83mm & 7,75mm & \textcolor{badred}{Kém -- không nên thay thế trạm thật} \\
\bottomrule
\end{tabular}
\vspace{0.3cm}

\small So ERA5 nội suy Bilinear vào tọa độ trạm với giá trị trạm thật đo được, cùng ngày
(108.542 cặp nhiệt độ, 61 trạm, 2020--2024). RMSE nhiệt độ theo tercile độ cao: 0,99 / 0,94 /
\textbf{1,42}°C -- nhóm trạm cao nhất sai số lớn nhất, đúng vật lý (lưới 0,25° không bắt được địa hình).
\end{frame}

\begin{frame}{Cache giúp truy vấn Spark SQL nhanh hơn 68\,\% thật trên 9,3 triệu bản ghi, vượt ngưỡng 30\,\% yêu cầu}
\begin{center}
\bignum{68,1\%}{cải thiện tốc độ\\truy vấn (cache)}{goodgreen}{5cm}
\bignumS{0,37s $\to$ 0,12s}{thời gian trung bình\\3 lần chạy}{hustblue}{5cm}
\end{center}
\vspace{0.2cm}
\centering\small
\texttt{warehouse/query\_optimization.py} -- EXPLAIN 2 bản (Scan parquet thô vs InMemoryTableScan)
cho truy vấn tổng hợp nhiệt độ theo (trạm, tháng), nguồn GHCN.
\end{frame}

\begin{frame}{Giảm data skew từ 10,6\,\% xuống 0,9\,\% khả thi nhưng đánh đổi 12 lần số partition -- quyết định giữ nguyên}
\centering\small
\begin{tabular}{@{}lrr@{}}
\toprule
 & \texttt{year/source} (hiện tại) & \texttt{year/month/source} (thử nghiệm) \\
\midrule
Số partition & 109 & 1.306 \\
Partition lớn nhất & 10,62\% & \textcolor{goodgreen}{0,90\%} \\
\bottomrule
\end{tabular}
\vspace{0.3cm}

\small \textbf{Quyết định trung thực}: KHÔNG áp dụng -- skew hiện tại (10,6\%) đã thấp, còn 1.306 partition
làm nặng thêm vấn đề ``small-file'': đo lại công bằng trên 9,3 triệu dòng, partition hiện tại vẫn
\textbf{chậm hơn 1,5--2,4 lần} so với không partition, vì Parquet đã tự bỏ qua dữ liệu không liên quan nhờ
thống kê min/max theo row-group. Lợi ích partition pruning chỉ rõ ở quy mô lớn hơn hoặc trên S3/GCS.
\end{frame}

%=============================================================
\section{Spark, Streaming \& PINN}
%=============================================================
\hypertarget{sec4}{}
\hustsectionpage{4. Spark, Streaming\\\& PINN (TV3)}

\begin{frame}{PINN trường gió không vượt baseline hồi quy thuần ở quy mô dữ liệu hiện tại -- một kết quả trung thực, không che giấu}
\centering\small
\begin{tabular}{@{}crrr@{}}
\toprule
$\lambda$ & Test RMSE & Test MAE & Physics loss cuối \\
\midrule
0,0 (baseline) & \textbf{1,0447} & 0,7429 & 0 \\
0,01 & 1,0889 & 0,7648 & 0,5303 \\
0,1 & 1,0826 & 0,7661 & 0,1192 \\
1,0 & 1,1227 & 0,7927 & 0,0119 \\
\bottomrule
\end{tabular}

\vspace{0.6cm}
\centering\small
Dữ liệu: 113 trạm ERA5 thật (gió trung bình 15--17/6/2021), ràng buộc continuity Navier-Stokes 2D rút gọn. Nguyên nhân hợp lý:
chỉ 1 snapshot không-thời-gian, ràng buộc vật lý cần nhiều điểm/chuỗi thời gian hơn để thể hiện
lợi thế MSE. Đã verify tái lập được tới 4 chữ số thập phân trên venv hoàn toàn sạch.
\end{frame}

\begin{frame}{Spark Structured Streaming đọc Kafka thật và tính đúng cửa sổ trượt 1 phút, phát hiện thêm 2 bug môi trường}
\begin{center}
\bignumS{\checkmark\ 20/10}{Spark 3 node + Kafka\\đã verify thật}{goodgreen}{4.5cm}
\bignumS{\checkmark\ 09/11}{Streaming windowing\\test live thành công}{goodgreen}{4.5cm}
\end{center}
\vspace{0.2cm}
\centering\small
\texttt{warehouse/spark\_streaming\_consumer.py} đọc topic Kafka thật, console hiển thị đúng
\texttt{avg\_temp\_c} theo cửa sổ trượt 1 phút/30s từng trạm. Phát hiện + sửa 2 bug môi trường
thật: (1) \texttt{SparkContext} lỗi do hostname tự phân giải ra IP LAN thay vì localhost (nghi do
mạng ảo Docker Desktop vừa khởi động lại) -- sửa bằng \texttt{SPARK\_LOCAL\_IP=127.0.0.1}; (2)
Kafka connector ghim cứng version 4.0.0 không khớp PySpark cài thật 4.2.0 -- sửa bằng lấy version
tự động.
\end{frame}

%=============================================================
\section{Phát hiện dị thường khí hậu}
%=============================================================
\hypertarget{sec4b}{}
\hustsectionpage{Phát hiện dị thường\\khí hậu}

\begin{frame}{Hệ thống phát hiện đúng cả 3 sự kiện khí hậu có thật -- nắng nóng 2024 nổi bật hẳn so với mọi năm}
\begin{columns}[T]
\begin{column}{0.56\textwidth}
\centering\scriptsize
Tỷ lệ ngày nóng bất thường ($z \ge 2$), tháng 4--5, 34 trạm GHCN\\[2pt]
\vspace{-0.15cm}
\begin{tikzpicture}[x=0.52cm, y=0.17cm]
\foreach \yr/\v in {2016/3.5,2017/0.2,2018/0.5,2019/3.9,2020/2.4,2021/0.8,2022/0.4,2023/4.1,2025/0.8} {
  \pgfmathsetmacro{\xx}{\yr-2016}
  \fill[gray!55, rounded corners=0.6pt] (\xx+0.12,0) rectangle (\xx+0.88,\v);
  \node[font=\tiny, rotate=90, anchor=east] at (\xx+0.5,-0.2) {\yr};
}
\fill[badred, rounded corners=0.6pt] (8.12,0) rectangle (8.88,11.2);
\node[font=\tiny, rotate=90, anchor=east] at (8.5,-0.2) {\textbf{2024}};
\node[font=\scriptsize, above, text=badred] at (8.5,11.2) {\textbf{11,2\%}};
\draw[dashed, hustblue] (0,1.2) -- (10,1.2) node[right, font=\tiny, align=left] {nền\\1991--2020\\1,2\%};
\draw[gray] (0,0) -- (10,0);
\end{tikzpicture}
\end{column}
\begin{column}{0.44\textwidth}
\scriptsize
\textbf{Nắng nóng 4--5/2024}: 12,3\% ngày $z \ge 2$, gấp \textbf{10 lần} cùng kỳ 1991--2020; 29/34 trạm có
đợt nóng; Thái Lan 30 ngày liên tục, TP.HCM 27 ngày.\\[4pt]
\textbf{Bão Yagi 7/9/2024}: ERA5 áp suất \textbf{978,8 hPa} tại Hải Phòng, $z = -7{,}2$ -- đúng nơi bão đổ bộ.\\[4pt]
\textbf{Hạn Mekong 2019--2020}: SPI-3 trung bình $-0{,}59$, tháng hạn 22\% (bình thường 13\%) -- tín hiệu yếu
do chỉ 2 trạm đủ số liệu mưa.\\[4pt]
\textit{Phương pháp chuẩn WMO/ETCCDI}: khí hậu nền theo ngày-trong-năm $\pm$15 ngày (GHCN 1991--2020),
đợt nóng TG90p, SPI khớp gamma -- tính bằng Spark trên 977 nghìn ngày-trạm.
\end{column}
\end{columns}
\end{frame}

\begin{frame}{Tích hợp trạm với ERA5 tách được dị thường khí hậu thật khỏi cảm biến hỏng -- điều 1 nguồn đơn lẻ không làm được}
\centering\small
\begin{tabular}{@{}lllr@{}}
\toprule
Trạm & ERA5 cùng vị trí, cùng ngày & Kết luận & Số trường hợp \\
\midrule
Bất thường & Cùng chiều, $|z| \ge 1$ & \textcolor{goodgreen}{Dị thường khí hậu thật} & 5.375 \\
Cực đoan ($|z| \ge 3$) & Bình thường hoặc ngược chiều & \textcolor{badred}{Nghi lỗi dữ liệu} & 25 \\
Bất thường & Lệch nhẹ & Dị thường cục bộ & 514 \\
\bottomrule
\end{tabular}
\vspace{0.3cm}

\small Ví dụ thật: trạm ISD \textbf{Phnom Penh} ghi áp suất \textbf{981,6 hPa ngày 3/7/2023} ($z = -9{,}1$) trong
khi ERA5 cùng vị trí bình thường ($z = +1{,}6$), và lặp lại suốt 7--9/2023 -- dấu hiệu cảm biến áp suất
hỏng/trôi, không phải bão. Dùng Record Linkage không-thời gian (ERA5 nội suy đúng vào \texttt{location\_id} của trạm).
\end{frame}

%=============================================================
\section{Dashboard}
%=============================================================
\hypertarget{sec5}{}
\hustsectionpage{5. Dashboard\\(TV4)}

\begin{frame}{Dashboard 7 panel đọc trực tiếp 9,3 triệu bản ghi thật, đã verify qua trình duyệt thật}
\centering\footnotesize
\begin{tabular}{@{}p{3.6cm}p{8.4cm}@{}}
\toprule
Panel & Nội dung thật \\
\midrule
1. Chuỗi thời gian theo trạm & 119 trạm/điểm, 6 biến chọn được \\
2. Tổng hợp theo nguồn & Bảng thống kê trực tiếp từ warehouse \\
3. Bản đồ + time slider & 113 trạm, join Location, mặc định ngày đủ dữ liệu \\
4. So sánh nguồn & NASA\_POWER vs OPEN\_METEO 2020--2024, RMSE thật \\
5. Cảnh báo bất thường & IQR thật, màu \textcolor{goodgreen}{xanh}/\textcolor{warnorange}{vàng}/\textcolor{badred}{đỏ} theo \% outlier \\
6. Output PINN & Bản đồ residual + bảng pred vs obs từng trạm \\
7. Dị thường khí hậu & Bản đồ $z$ (xanh--xám--đỏ), đợt nóng, SPI-3, xác nhận chéo ERA5 \\
\bottomrule
\end{tabular}

\vspace{0.4cm}
\centering\small
Đã verify thật bằng trình duyệt: cả 7 panel render đúng, 0 lỗi. Có cache bảng theo ngày: đổi ngày/trạm \textasciitilde 0,6\,s.
\end{frame}

%=============================================================
\section{Rủi ro, giới hạn \& kết luận}
%=============================================================
\hypertarget{sec6}{}
\hustsectionpage{6. Rủi ro, giới hạn\\\& kết luận}

\begin{frame}{Các giới hạn kỹ thuật còn tồn đọng đều đã ghi nhận rõ nguyên nhân, không che giấu}
\centering\scriptsize
\begin{tabular}{@{}p{4.2cm}p{7.9cm}@{}}
\toprule
Giới hạn & Nguyên nhân / hướng xử lý \\
\midrule
Khí hậu nền ngắn với ERA5/ISD (5--10 năm) & Kết luận khí hậu chỉ dựa trên GHCN nền WMO 1991--2020; SPI mới có 12 trạm do GHCN thiếu số liệu mưa \\
Pipeline inference (\texttt{infer.py}) chưa tách checkpoint & Script hiện train+eval trong 1 lần chạy -- đủ cho quy mô bài tập lớn, chưa tách để demo nhanh \\
PINN chưa vượt baseline RMSE thuần & Thiếu chiều thời gian trong dữ liệu hiện tại -- physics loss vẫn có giá trị về tính hợp lý vật lý \\
GitHub Projects board chưa thiết lập & Đã dùng Issues/Pull Request trên GitHub; trạng thái tác vụ theo dõi ở \texttt{docs/03} \\
\bottomrule
\end{tabular}
\end{frame}

\begin{frame}{120/160 tác vụ đã rà soát với bằng chứng thật, pipeline đầy đủ từ ingest tới dashboard}
\begin{center}
\bignum{79/120}{tác vụ Xong đúng\\tiêu chí kiểm thử}{goodgreen}{4.2cm}
\bignumS{9.348.470}{bản ghi thật\\5 nguồn, 9 nước}{goodgreen}{4.2cm}
\bignum{54}{test tự động\\(gồm Spark + Kafka thật)}{goodgreen}{4.2cm}
\end{center}
\vspace{0.4cm}
\centering\small
Toàn bộ số liệu trong báo cáo này tái lập được bằng cách chạy lại script tương ứng (xem README
mục 6b, \texttt{requirements.txt}, \texttt{CHANGELOG.md}) -- không có số liệu ước lượng hoặc bịa.
\end{frame}

\hustthankyou

\end{document}
